\section*{Sets and Indices}

\begin{itemize}
    \item Product models $m \in M = \{A,B\}$.
    \item Processes $p \in P = \{\mathrm{I},\mathrm{II}\}$.
\end{itemize}

\section*{Parameters}

All parameter values below are taken from the CSV files used by the implementation.

\begin{itemize}
    \item Process-I hours required per unit:
    \[
    a^{\mathrm{I}}_A = 4, \qquad a^{\mathrm{I}}_B = 6.
    \]
    \item Process-II hours required per unit:
    \[
    a^{\mathrm{II}}_A = 3, \qquad a^{\mathrm{II}}_B = 2.
    \]
    \item Regular weekly Process-II capacity:
    \[
    C^{\mathrm{II}} = 70.
    \]
    \item Required exact weekly Process-I usage:
    \[
    H^{\mathrm{I}} = 150.
    \]
    \item Maximum allowable Process-II overtime:
    \[
    O^{\max} = 30.
    \]
    \item Unit profit under regular production:
    \[
    \pi_A = 300, \qquad \pi_B = 450.
    \]
    \item Profit reduction for units produced in Process-II overtime:
    \[
    \delta_A = 20, \qquad \delta_B = 25.
    \]
    Hence overtime-unit profits are
    \[
    \pi^{o}_A = \pi_A - \delta_A = 280, \qquad
    \pi^{o}_B = \pi_B - \delta_B = 425.
    \]
    \item Minimum weekly profit requirement:
    \[
    \Pi^{\min} = 10000.
    \]
    \item Minimum production requirements:
    \[
    L_A = 10, \qquad L_B = 15.
    \]
\end{itemize}

\section*{Decision Variables}

\begin{itemize}
    \item $x^r_A$ (code: \texttt{xA\_r}): units of model A produced using regular Process-II capacity.
    \item $x^o_A$ (code: \texttt{xA\_o}): units of model A produced using overtime Process-II capacity.
    \item $x^r_B$ (code: \texttt{xB\_r}): units of model B produced using regular Process-II capacity.
    \item $x^o_B$ (code: \texttt{xB\_o}): units of model B produced using overtime Process-II capacity.
\end{itemize}

All variables are continuous and nonnegative:
\[
x^r_A, x^o_A, x^r_B, x^o_B \ge 0.
\]

For convenience, total production is
\[
X_A = x^r_A + x^o_A, \qquad X_B = x^r_B + x^o_B,
\]
corresponding in code to \texttt{total\_A} and \texttt{total\_B}.

\section*{Objective}

Maximize total weekly profit:
\[
\max \; Z
=
\pi_A x^r_A + \pi_B x^r_B + \pi^o_A x^o_A + \pi^o_B x^o_B,
\]
i.e.,
\[
\max \; Z
=
300x^r_A + 450x^r_B + 280x^o_A + 425x^o_B.
\]

\section*{Constraints}

\begin{align}
a^{\mathrm{I}}_A (x^r_A + x^o_A) + a^{\mathrm{I}}_B (x^r_B + x^o_B) &= H^{\mathrm{I}}
&& \text{(exact Process-I usage)} \\
a^{\mathrm{II}}_A x^r_A + a^{\mathrm{II}}_B x^r_B &\le C^{\mathrm{II}}
&& \text{(regular Process-II capacity)} \\
a^{\mathrm{II}}_A x^o_A + a^{\mathrm{II}}_B x^o_B &\le O^{\max}
&& \text{(Process-II overtime limit)} \\
x^r_A + x^o_A &\ge L_A
&& \text{(minimum production of A)} \\
x^r_B + x^o_B &\ge L_B
&& \text{(minimum production of B)} \\
\pi_A x^r_A + \pi_B x^r_B + \pi^o_A x^o_A + \pi^o_B x^o_B &\ge \Pi^{\min}
&& \text{(minimum weekly profit).}
\end{align}

Substituting the numerical values, the implemented model is
\begin{align}
4(x^r_A + x^o_A) + 6(x^r_B + x^o_B) &= 150, \\
3x^r_A + 2x^r_B &\le 70, \\
3x^o_A + 2x^o_B &\le 30, \\
x^r_A + x^o_A &\ge 10, \\
x^r_B + x^o_B &\ge 15, \\
300x^r_A + 450x^r_B + 280x^o_A + 425x^o_B &\ge 10000.
\end{align}

\section*{Notes on Implementation}

\begin{itemize}
    \item The source code first sketches an alternative nonlinear overtime formulation using total production and an overtime variable \texttt{ot}, but that model is not solved.
    \item The implemented and solved model is the second formulation, which splits each product into regular-capacity and overtime-capacity production for Process II.
    \item Overtime is modeled implicitly through the variables $x^o_A$ and $x^o_B$ and the constraint
    \[
    3x^o_A + 2x^o_B \le 30,
    \]
    rather than through a separate overtime-hours decision variable in the final solved model.
    \item Process I is enforced as an equality, exactly matching \texttt{process\_I\_hours}, not merely as an upper bound.
    \item The variables are continuous in the implementation; no integrality restrictions are imposed.
    \item The minimum-profit requirement is included explicitly as a constraint even though the same expression is maximized in the objective.
    \item The model is solved as a linear program via PuLP with the GUROBI command interface (\texttt{pulp.GUROBI\_CMD}).
\end{itemize}
